\documentclass[11pt]{article}
\usepackage[margin=1.1in]{geometry}
\usepackage{amsmath,amssymb,amsthm}
\usepackage{hyperref}

\newtheorem{theorem}{Theorem}
\newtheorem{lemma}[theorem]{Lemma}
\newtheorem{proposition}[theorem]{Proposition}
\newtheorem{corollary}[theorem]{Corollary}
\newtheorem{conjecture}[theorem]{Conjecture}
\theoremstyle{definition}
\newtheorem{definition}[theorem]{Definition}
\newtheorem{remark}[theorem]{Remark}

\newcommand{\Ecal}{\mathcal{E}}
\newcommand{\Xcal}{\mathcal{X}}
\newcommand{\Scal}{\mathcal{S}}
\newcommand{\Fcal}{\mathcal{F}}

\title{Ergodicity of the projection-preserving switch chain\\
on deduplicated single-sender directed hypergraphs}
\author{}
\date{}

\begin{document}
\maketitle

\begin{abstract}
We record the state space, the move set, and the ergodic properties of the
Metropolis chain used to sample the exponential tilt $P_\theta(H)\propto
e^{\theta\Phi(H)}$ over regroupings of a fixed directed projection. Symmetry of
the proposal, aperiodicity, and the reduction to a per-sender problem are proved
in full. Irreducibility is proved on the \emph{unrestricted} (row-repetition
permitted) state space by reduction to Ryser's interchange theorem, and on the
\emph{deduplicated} state space in the case where every hyperedge has size two,
by reduction to the classical switch-connectivity theorem for simple graphs. For
the deduplicated space with general size profiles we state the connectivity
claim as a conjecture and report exhaustive computational verification; we are
not aware of a proof, and we note that the analogous gap is present in the
literature this construction builds on.
\end{abstract}

\section{Setup}

\begin{definition}[Single-sender directed hypergraph]
Let $V$ be a finite node set. An \emph{event} is a pair $(s,R)$ with $s\in V$
and $R\subseteq V\setminus\{s\}$, $R\neq\emptyset$; $s$ is the \emph{sender}
(tail) and $R$ the \emph{recipient set} (head). The \emph{team} of the event is
$N=\{s\}\cup R$ and its \emph{size} is $k=|N|$. A single-sender directed
hypergraph is a finite collection $H$ of events.
\end{definition}

Following Abuissa et al.~\cite{dinghy}, we distinguish \emph{edge-ordered}
hypergraphs, in which hyperedges carry unique identifiers, from
\emph{edge-unordered} ones, in which they do not. We work throughout with the
\emph{deduplicated} (or \emph{simple}) case, in which $H$ is a \emph{set} of
events: no event is repeated. This is strictly stronger than being
edge-unordered, where the edge collection is a multiset whose elements merely
lack identifiers.

\begin{definition}[Projection and fiber]
The \emph{weighted directed projection} of $H$ is the map
$\pi_H:V\times V\to\mathbb{Z}_{\ge0}$,
\[
  \pi_H(u,v)\;=\;\#\{(s,R)\in H:\ s=u,\ v\in R\}.
\]
Given an observed $H_0$, the \emph{fiber} $\Fcal(H_0)$ is the set of
deduplicated single-sender directed hypergraphs $H$ such that
\begin{enumerate}
\item[(F1)] $\pi_H=\pi_{H_0}$ (identical weighted directed projection), and
\item[(F2)] for every sender $s$, the multiset of sizes of $s$'s events in $H$
equals that in $H_0$.
\end{enumerate}
\end{definition}

Condition (F1) is what makes graph reciprocity constant across the fiber;
(F2) is an identification choice, conditioning out the hyperedge-size
composition, which mechanically influences every group-reciprocity measure.

\begin{definition}[Switch]
Let $H\in\Fcal(H_0)$, let $s$ be a sender, and let $E_i\neq E_j$ be two of $s$'s
recipient sets. Let $a\in E_i\setminus E_j$ and $b\in E_j\setminus E_i$. The
\emph{switch} $\sigma_{s,i,j,a,b}$ replaces $E_i$ by
$E_i'=(E_i\setminus\{a\})\cup\{b\}$ and $E_j$ by
$E_j'=(E_j\setminus\{b\})\cup\{a\}$, leaving all other events unchanged. The
switch is \emph{admissible} if the resulting collection is still deduplicated.
\end{definition}

\begin{definition}[The chain]
Fix $\theta\in\mathbb{R}$ and a statistic $\Phi:\Fcal(H_0)\to\mathbb{R}$. One
step of the chain from $H$: choose a sender $s$ uniformly from those with at
least two events, choose ordered indices $i,j$ uniformly and independently from
$s$'s $m_s$ events, choose $a$ uniformly from $E_i\setminus E_j$ and $b$
uniformly from $E_j\setminus E_i$. If $i=j$, or either difference set is empty,
or the switch is inadmissible, remain at $H$. Otherwise move to
$H'=\sigma_{s,i,j,a,b}(H)$ with probability $\min\{1,e^{\theta(\Phi(H')-\Phi(H))}\}$
and remain at $H$ otherwise.
\end{definition}

\begin{remark}
Every failed proposal consumes a step and is a self-loop. Resampling until a
valid proposal is found would replace the transition kernel $p(H,\cdot)$ by
$p(H,\cdot)/(1-p_{\mathrm{fail}}(H))$; since $p_{\mathrm{fail}}$ is
state-dependent, the resulting chain is not symmetric and its stationary law is
proportional to $e^{\theta\Phi(H)}\bigl(1-p_{\mathrm{fail}}(H)\bigr)$ rather
than $e^{\theta\Phi(H)}$.
\end{remark}

\section{Structural lemmas}

\begin{lemma}[Invariance]\label{lem:inv}
Every switch preserves (F1) and (F2). Consequently the chain never leaves
$\Fcal(H_0)$.
\end{lemma}

\begin{proof}
A switch moves recipient $a$ from $E_i$ to $E_j$ and $b$ from $E_j$ to $E_i$,
both within the events of a single sender $s$. For $r\notin\{a,b\}$ the count
$\pi(s,r)$ is untouched; $\pi(s,a)$ loses one occurrence in $E_i$ and gains one
in $E_j$, and symmetrically for $b$; no other sender's events change. Hence (F1).
Since $|E_i'|=|E_i|$ and $|E_j'|=|E_j|$, the multiset of sizes is preserved,
giving (F2). Admissibility preserves deduplication by definition.
\end{proof}

\begin{lemma}[Involution and symmetry]\label{lem:sym}
Write $A=E_i\setminus E_j$, $B=E_j\setminus E_i$, $C=E_i\cap E_j$. After the
switch, $A'=E_i'\setminus E_j'=(A\setminus\{a\})\cup\{b\}$ and
$B'=E_j'\setminus E_i'=(B\setminus\{b\})\cup\{a\}$; in particular $|A'|=|A|$ and
$|B'|=|B|$. The map $\sigma_{s,i,j,b,a}$ inverts $\sigma_{s,i,j,a,b}$, and the
proposal kernel $q$ satisfies $q(H,H')=q(H',H)$.
\end{lemma}

\begin{proof}
$E_i'=C\cup(A\setminus\{a\})\cup\{b\}$ and $E_j'=C\cup(B\setminus\{b\})\cup\{a\}$.
As $A,B,C$ are pairwise disjoint and $a\neq b$, the stated forms of $A'$ and
$B'$ follow, and the cardinalities are preserved. Applying the switch with the
roles of $a$ and $b$ exchanged restores $E_i$ and $E_j$. The proposal
probability of the forward move is
$\frac{1}{|S|}\cdot\frac{1}{m_s^2}\cdot\frac{1}{|A|}\cdot\frac{1}{|B|}$
and of the reverse move
$\frac{1}{|S|}\cdot\frac{1}{m_s^2}\cdot\frac{1}{|A'|}\cdot\frac{1}{|B'|}$;
these agree.
\end{proof}

Lemma~\ref{lem:sym} is the analogue in our setting of \cite[Theorem~2]{dinghy},
which establishes double stochasticity of the DiNgHy transition matrix.

\begin{lemma}[Identity switches]\label{lem:id}
$|A|=|B|=1$ if and only if the switch exchanges $E_i$ and $E_j$, i.e.\
$E_i'=E_j$ and $E_j'=E_i$. In the deduplicated (edge-unordered) state space such
a switch is the identity.
\end{lemma}

\begin{proof}
If $A=\{a\}$ and $B=\{b\}$ then $E_i=C\cup\{a\}$ and $E_j=C\cup\{b\}$, so
$E_i'=C\cup\{b\}=E_j$ and $E_j'=C\cup\{a\}=E_i$. Conversely if $E_i'=E_j$ then
$(A\setminus\{a\})\cup\{b\}=B$, and since $b\notin A$ this forces $A=\{a\}$;
symmetrically $B=\{b\}$.
\end{proof}

\begin{lemma}[Factorisation]\label{lem:fact}
$\Fcal(H_0)\cong\prod_{s}\Fcal_s(H_0)$, where $\Fcal_s$ is the set of
deduplicated collections of $s$'s events satisfying (F1) and (F2) for that
sender alone. Every switch acts within a single factor. Hence the chain is
irreducible on $\Fcal(H_0)$ if and only if its restriction to $\Fcal_s$ is
irreducible for every $s$.
\end{lemma}

\begin{proof}
Conditions (F1) and (F2) are stated per sender, and by Lemma~\ref{lem:inv} a
switch modifies only the events of one sender. Deduplication also decomposes:
two events with different senders are automatically distinct. The product
structure and the componentwise action follow, and irreducibility of a product
chain in which each step moves one coordinate is equivalent to irreducibility of
each coordinate chain.
\end{proof}

\medskip
By Lemma~\ref{lem:fact} we may fix a single sender $s$ and drop it from the
notation. Let $R_1,\dots,R_n$ enumerate the recipients ever contacted by $s$,
let $c_r=\pi(s,R_r)$, and let $E_1,\dots,E_m$ be $s$'s events. Encode the state
as the $m\times n$ incidence matrix $M$ with $M_{ir}=\mathbf{1}[R_r\in E_i]$.
Then (F2) fixes the row sums, (F1) fixes the column sums, and deduplication
says the rows of $M$ are pairwise distinct. A switch is exactly the classical
\emph{interchange}: choosing rows $i\neq j$ and columns $a\neq b$ with
$M_{ia}=M_{jb}=1$, $M_{ib}=M_{ja}=0$, and replacing the $2\times2$ submatrix
$\begin{psmallmatrix}1&0\\0&1\end{psmallmatrix}$ by
$\begin{psmallmatrix}0&1\\1&0\end{psmallmatrix}$.

\begin{lemma}[Aperiodicity]\label{lem:aper}
For every $H\in\Fcal(H_0)$, $p(H,H)\ge c$ where
$c=\frac{1}{|S|}\sum_{s\in S}\frac{1}{m_s}>0$ is independent of $H$. In
particular the chain is aperiodic.
\end{lemma}

\begin{proof}
The event $\{i=j\}$ has probability $1/m_s$ once sender $s$ is chosen, and
results in a self-loop. By (F2) the number $m_s$ of events of each sender is
invariant along the chain, so $c$ does not depend on the state.
\end{proof}

\begin{remark}
Lemma~\ref{lem:aper} shows that making the chain \emph{lazy} --- replacing $p$
by $\lambda p+(1-\lambda)I$ --- is unnecessary here: the required uniform lower
bound on the holding probability is already available. Laziness preserves the
stationary law and reversibility, but multiplies the mixing time by
$\Theta(1/\lambda)$, and \cite{dinghy} indeed report designing their proposal to
\emph{reduce} self-loops for faster mixing.
\end{remark}

\section{Irreducibility}

\subsection{The unrestricted state space}

Let $\Xcal$ denote the set of $m\times n$ zero-one matrices with the given row
and column sums, \emph{without} the distinct-rows requirement --- equivalently,
the edge-ordered state space in which a sender may repeat an event.

\begin{theorem}[Ryser]\label{thm:ryser}
Any two matrices in $\Xcal$ are connected by a finite sequence of interchanges.
Consequently the switch chain restricted to $\Xcal$ is irreducible, and together
with Lemmas~\ref{lem:sym} and \ref{lem:aper} it is ergodic with stationary law
proportional to $e^{\theta\Phi}$.
\end{theorem}

\begin{proof}
This is Ryser's interchange theorem \cite{ryser1957}; see also Brualdi's survey
\cite{brualdi1980}. In the bipartite-graph formulation --- rows and columns as
the two sides, $M_{ir}=1$ as an edge --- it is the statement that the switch
chain on bipartite graphs with prescribed degree sequences is irreducible, as
used by Kannan, Tetali and Vempala \cite{ktv1999} and by
Abuissa et al.~\cite{dinghy}.
\end{proof}

\subsection{The deduplicated state space}

Let $\Scal\subsetneq\Xcal$ be the subset of matrices with pairwise distinct
rows, and let switches be admissible only when they stay in $\Scal$. Removing
states from a connected graph can disconnect it, so Theorem~\ref{thm:ryser}
does \emph{not} imply that $\Scal$ is connected, and the question is genuinely
different.

Two special cases are settled.

\begin{proposition}[Distinct sizes]\label{prop:sizes}
If the row sums $|E_1|,\dots,|E_m|$ are pairwise distinct then $\Scal=\Xcal$ and
the chain is irreducible on $\Scal$ by Theorem~\ref{thm:ryser}.
\end{proposition}

\begin{proof}
Rows of different sums cannot be equal, so the distinctness constraint is
vacuous, and every interchange is admissible.
\end{proof}

\begin{theorem}[All events dyadic]\label{thm:size2}
If every event has size $k=3$, i.e.\ every row sum equals $2$, then the switch
chain is irreducible on $\Scal$.
\end{theorem}

\begin{proof}
Identify each row, a $2$-subset of recipients, with an edge of a graph $G$ on
the recipient vertex set; the column sums are the degrees of $G$ and
distinctness of rows says exactly that $G$ is simple. An admissible interchange
takes two disjoint or sharing edges and rewires them without creating a
multi-edge, which is precisely the double edge swap. The theorem is then the
classical statement that any two simple graphs with the same degree sequence are
connected by double edge swaps, due to Petersen \cite{petersen1891} in the
regular case and in the general form given by Havel \cite{havel1955}, Hakimi
\cite{hakimi1962}, Eggleton and Holton \cite{eggleton1979} and Taylor
\cite{taylor1981}. See \cite{fosdick2018} for a modern treatment aimed at
sampling.
\end{proof}

For general size profiles we do not have a proof.

\begin{conjecture}\label{conj:main}
For every choice of row sums and column sums, the admissible-interchange graph
on $\Scal$ is connected.
\end{conjecture}

\paragraph{Computational verification.}
Conjecture~\ref{conj:main} was checked exhaustively on small instances. For each
$n\le 7$ recipients and each size profile drawn from uniform profiles
$(k,\dots,k)$ with $k\in\{2,3,4\}$ and $m\le 6$ rows, and mixed profiles
combining sizes $2,3,4$, every fiber was enumerated, grouped by column-sum
vector, and the admissible-interchange graph on each fiber with at least two
states was tested for connectivity by breadth-first search. Across two
independent sweeps this covered roughly $4.3\times10^4$ non-trivial fibers, the
largest containing $2367$ states. \emph{No disconnected fiber was found.}

\begin{remark}
This is evidence, not proof, and the reader should treat
Conjecture~\ref{conj:main} accordingly. We note that the same gap is present in
the closest prior work: Abuissa et al.~\cite{dinghy} show that the Markov chain
underlying their sampler is \emph{not} irreducible in general, and supply only
sufficient conditions for irreducibility; they also observe that no general
mixing-time analysis is available even for the dyadic switch chain
\cite{erdos2022}. A practical mitigation, which we adopt, is to verify
empirically that chains started from several distinct states in the same fiber
converge to a common distribution of $\Phi$.
\end{remark}

\section{Consequence}

\begin{corollary}
Assume Conjecture~\ref{conj:main}, or either of the hypotheses of
Proposition~\ref{prop:sizes} and Theorem~\ref{thm:size2}. Then the chain is
irreducible (Lemma~\ref{lem:fact}) and aperiodic (Lemma~\ref{lem:aper}) on
$\Fcal(H_0)$, its proposal is symmetric (Lemma~\ref{lem:sym}), and hence the
Metropolis chain has unique stationary distribution
\[
  P_\theta(H)\;=\;\frac{e^{\theta\Phi(H)}}{\sum_{H'\in\Fcal(H_0)}e^{\theta\Phi(H')}},
  \qquad H\in\Fcal(H_0).
\]
At $\theta=0$ this is the uniform distribution on $\Fcal(H_0)$, i.e.\ the
projection-preserving null model. Since every $H\in\Fcal(H_0)$ has
$\pi_H=\pi_{H_0}$, the reciprocity of the directed projection is constant on
$\Fcal(H_0)$ and therefore invariant in $\theta$.
\end{corollary}

\begin{remark}
Because the state space is deduplicated, $|E|=\sum_N p_N$, the number of
distinct directed hyperedges, is constant on $\Fcal(H_0)$. As $|E|$ is the
denominator of the group-reciprocity measures, those measures are then computed
against a fixed normalisation for every $\theta$. This fails on the
edge-ordered space, where an interchange may render two of a sender's events
identical.
\end{remark}

\begin{thebibliography}{99}

\bibitem{dinghy}
M.~Abuissa, A.~Kirkley, and M.~Riondato.
\newblock DiNgHy: null models for non-degenerate directed hypergraphs.
\newblock \emph{Data Mining and Knowledge Discovery} (to appear).
\newblock \emph{[Verify authors, year, volume before use.]}

\bibitem{ryser1957}
H.~J.~Ryser.
\newblock Combinatorial properties of matrices of zeros and ones.
\newblock \emph{Canadian Journal of Mathematics}, 9:371--377, 1957.

\bibitem{brualdi1980}
R.~A.~Brualdi.
\newblock Matrices of zeros and ones with fixed row and column sum vectors.
\newblock \emph{Linear Algebra and its Applications}, 33:159--231, 1980.

\bibitem{ktv1999}
R.~Kannan, P.~Tetali, and S.~Vempala.
\newblock Simple Markov-chain algorithms for generating bipartite graphs and
  tournaments.
\newblock \emph{Random Structures \& Algorithms}, 14(4):293--308, 1999.

\bibitem{petersen1891}
J.~Petersen.
\newblock Die Theorie der regul\"aren Graphs.
\newblock \emph{Acta Mathematica}, 15:193--220, 1891.

\bibitem{havel1955}
V.~Havel.
\newblock A remark on the existence of finite graphs.
\newblock \emph{\v{C}asopis pro p\v{e}stov\'an\'i matematiky}, 80:477--480, 1955.

\bibitem{hakimi1962}
S.~L.~Hakimi.
\newblock On realizability of a set of integers as degrees of the vertices of a
  linear graph.
\newblock \emph{SIAM Journal on Applied Mathematics}, 10(3):496--506, 1962.

\bibitem{eggleton1979}
R.~B.~Eggleton and D.~A.~Holton.
\newblock Simple and multigraphic realizations of degree sequences.
\newblock In \emph{Combinatorial Mathematics VIII}, Lecture Notes in
  Mathematics 884, Springer, 1979.

\bibitem{taylor1981}
R.~Taylor.
\newblock Constrained switchings in graphs.
\newblock In \emph{Combinatorial Mathematics VIII}, Lecture Notes in
  Mathematics 884, pages 314--336. Springer, 1981.

\bibitem{fosdick2018}
B.~K.~Fosdick, D.~B.~Larremore, J.~Nishimura, and J.~Ugander.
\newblock Configuring random graph models with fixed degree sequences.
\newblock \emph{SIAM Review}, 60(2):315--355, 2018.

\bibitem{erdos2022}
P.~L.~Erd\H{o}s, T.~R.~Mezei, I.~Mikl\'os, et al.
\newblock On the mixing time of the switch Markov chain.
\newblock \emph{[Cited via \cite{dinghy}; verify exact reference.]}

\end{thebibliography}

\end{document}
